\documentclass[10.5pt]{article}
\usepackage{amsmath, amsfonts, amssymb,amsthm}
\usepackage[includeheadfoot]{geometry} % For page dimensions
\usepackage{fancyhdr}
\usepackage{enumerate} % For custom lists
\usepackage{tikz-cd}
\usepackage{graphicx}

\fancyhf{}
\lhead{Engel's Theorem}
\rhead{Tighe McAsey}
\pagestyle{fancy}

% Page dimensions
\geometry{a4paper, margin=1in}

\theoremstyle{definition}
\newtheorem{pb}{}
\newtheorem{definition}{Definition}
\newtheorem{proposition}{Proposition}
\newtheorem{remark}{Remark}
\newtheorem{example}{Example}
\newtheorem{theorem}{Theorem}
\numberwithin{example}{subsection}
\numberwithin{pb}{subsection}
\usepackage{tikz-cd, stackengine}

% Commands:

\newcommand{\set}[1]{\left\{#1\right\}}
\newcommand{\gen}[1]{\langle#1\rangle}
\newcommand{\abs}[1]{\left\vert#1\right\vert}
\newcommand{\norm}[1]{\lvert\lvert#1\rvert\rvert}
\newcommand{\tand}{\text{ and }}
\newcommand{\tor}{\text{ or }}
\newcommand{\pd}{\frac{\partial}{\partial x_j}}
\newcommand{\px}{\frac{\partial}{\partial x}}
\newcommand{\py}{\frac{\partial}{\partial y}}
\newcommand{\pz}{\frac{\partial}{\partial z}}
\newcommand{\ppx}{\frac{\partial^2}{\partial x^2}}
\newcommand{\ppy}{\frac{\partial^2}{\partial y^2}}
\newcommand{\ppz}{\frac{\partial^2}{\partial z^2}}
\newcommand{\hess}{\operatorname{Hess}}
\newcommand{\inv}{\operatorname{inv}}
\newcommand{\spl}{\operatorname{\mathfrak{sl}}}
\newcommand{\spg}{\mathfrak{g}}
\newcommand{\sph}{\mathfrak{h}}
\newcommand{\spu}{\operatorname{\mathfrak{su}}}
\newcommand{\spo}{\operatorname{\mathfrak{so}}}

\begin{document}
     \begin{proposition}\label{engelstheoremlemma}
            A subalgebra \(\mathfrak{g} \subset \mathfrak{gl}(V)\) of nilpotent endomorphisms has a common eigenvector.
            \begin{proof}
                Induct on the dimension of the lie algebra. Dimension \(1\) is trivial, since for \(0 \neq \xi \in \mathfrak{g}\), nilpotence of \(\xi\) implies all of its eigenvalues are zero and thus in \(k\) implying that \(\xi\) has a nonzero eigenvector.

                Proceeding by induction let \(\mathfrak{g} \subset \mathfrak{gl}(V)\) be a lie algebra of nilpotent endomorphisms having finite dimension greater than \(1\), and suppose the result holds for all lie algebras of dimension less that \(\dim \mathfrak{g}\). Let \(\mathfrak{h} \subset \mathfrak{g}\) be a nonzero proper subalgebra (one always exists since we can take the abelian subalgebra generated by a single element). Then we consider the vectorspace \(U = \mathfrak{g}/\mathfrak{h}\), the adjoint representation of \(\mathfrak{h}\) on \(U\) is nilpotent, thus by induction (since \(\mathfrak{h} \subsetneq \mathfrak{g}\) has lesser dimension), \(\mathfrak{h}\) acting on \(U\) has a common eigenvector. This shows that \(\mathfrak{h}\) is properly contained in its normalizer \(N_\mathfrak{g}(\mathfrak{h}) \supsetneq \mathfrak{h}\).

                Now that we have established that for \(\mathfrak{h} \subsetneq \mathfrak{g}\) a subalgebra that \(N_\mathfrak{g}(\mathfrak{h}) \supsetneq \mathfrak{h}\) we can see that \(\mathfrak{g}\) has a codimension 1 ideal. To see this, note that \(N_\mathfrak{g}(\mathfrak{h})\) is a subalgebra of \(\mathfrak{g}\) having greater dimension that \(\mathfrak{h}\). Denote
                \begin{align}
                    N^{(k)}_{\mathfrak{g}}(\mathfrak{h}) = \begin{cases}
                        \mathfrak{h} & k = 0 \\
                        N_\mathfrak{g}(N^{(k-1)}_{\mathfrak{g}}(\mathfrak{h})) & k > 0
                    \end{cases}
                \end{align}
                by the previous argumentation, for \(k > 0\) either \(N^{(k)}_\mathfrak{g}(\mathfrak{h}) = \mathfrak{g}\) or \(\dim N^{(k)}_\mathfrak{g}(\mathfrak{h}) > \dim N^{(k-1)}_\mathfrak{g}(\mathfrak{h})\), since each of these is a subalgebra of \(\mathfrak{g}\) and \(\dim \mathfrak{g}\) is finite the latter case can only happen finitely many times. Thus for some \(k \geq 0\), we have \(N_\mathfrak{g}^{(k)}(\mathfrak{h}) \subsetneq N_\mathfrak{g}^{(k+1)}(\mathfrak{h}) = \mathfrak{g}\). This establishes that \(\mathfrak{g}\) has a nonzero proper ideal, which I will denote as \(I\). To scaffold this into a codimension-\(1\) ideal I study the adjoint representation on \(\mathfrak{g}/I\). If \(\ker \operatorname{ad} = \mathfrak{g}/I\), then \(\mathfrak{g}/I\) is abelian, thus has a codimension \(1\) ideal and the preimage of this under the quotient map \(\mathfrak{g} \to \mathfrak{g}/I\) is a codimension \(1\) ideal of \(\mathfrak{g}\). If \(\ker \operatorname{ad} = 0\), then by the first isomorphism theorem \(\mathfrak{g}/I \cong \operatorname{ad}(\mathfrak{g}/I) \subset \mathfrak{gl}(\mathfrak{g}/I)\), then following the same argumentation used to produce \(I\) \textcolor{red}{(\(*\))}, we find that \(\operatorname{ad}(\mathfrak{g}/I)\) has a nonzero proper ideal \(J\), and thus by the correspondence theorem, \(I\) is contained in a larger proper ideal \(I \subsetneq J' \subsetneq \mathfrak{g}\), since \(\mathfrak{g}\) is finite dimensional iterating this process we eventually get a codimension \(1\) ideal. Finally, if \(0 \subsetneq \ker \operatorname{ad} \subsetneq \mathfrak{g}/I\), then by the correspondence theorem, \(\ker \operatorname{ad}\) corresponds to a proper ideal \(J'\) with \(I \subsetneq J' \subsetneq \mathfrak{g}\) so iterating this process we eventually get a codimension \(1\) ideal.

                We have established existence of a codimension \(1\) ideal \(\mathfrak{h} \subset \mathfrak{g}\). Applying the induction hypothesis, i.e. let \(W \subset V\) be the nonzero subspace of eigenvectors for \(\mathfrak{h}\), since \(\mathfrak{h}\) is an ideal we have for \(w \in W\), \(\xi \in \mathfrak{g}, \eta \in \mathfrak{h}\)
                \begin{align}
                    \eta\xi(w) = \xi \underbrace{\eta (w)}_{=0} + \underbrace{[\eta,\xi]}_{\in \mathfrak{h}}(w) = 0
                \end{align}
                So that \(\mathfrak{g}: W \to W\), now writing \(\mathfrak{g} = \mathfrak{h} + k\xi\), we know the nilpotent endomorphism \(\xi: W \to W\) has an eigenvector,  this is already a common eigenvector for \(\mathfrak{h}\) so it is a common eigenvector for \(k\xi \oplus \mathfrak{h} = \mathfrak{g}\).
            \end{proof}
        \end{proposition}

        \textcolor{red}{(\(*\))} To justify why the same argument works here, we first note that because every element of \(\mathfrak{g}\) is a nilpotent operator, every element is \(\operatorname{ad}\)-nilpotent (to see this simply compute the word \(\operatorname{ad}_x^{n}(y)\)), this of course implies \(\operatorname{ad}\)-nilpotence for \(\mathfrak{g}/I\), equivalently every element of \(\operatorname{ad}(\mathfrak{g}/I)\) is a nilpotent operator in \(\mathfrak{gl}(\mathfrak{g}/I)\). Moreover \(\dim \mathfrak{g}/I < \dim \mathfrak{g}\) so it is valid to apply the induction hypothesis.
        \newpage

        \begin{theorem}[Engel]\label{engel2}
            If \(\mathfrak{g} \subset \mathfrak{gl}(V)\) is a subalgebra of nilpotent operators, then there is a basis for \(V\) so that for every \(x \in \mathfrak{g}\) the matrix representing \(x\) is strictly upper triangular.

            \begin{proof}
                It suffices to show that there exists a sequence of subspaces \(V_j \subset V\) with \(\dim V_j = j\) such that
                \begin{align}
                    0 = V_0 \subset V_1 \subset \cdots \subset V_{n-1} \subset V_n = V
                \end{align}
                and \(\mathfrak{g}(V_j) \subset V_{j-1}\).

                By \textbf{Proposition \ref{engelstheoremlemma}} \(\mathfrak{g}\) has a common eigenvector \(v \in V\), its eigenvalue must be zero since \(\mathfrak{g}\) consists of nilpotent operators. Choose \(V_1 = \operatorname{span}_k\set{v}\). Now assume that \(V_1\subset\cdots\subset V_r\) are vectorspaces with \(\dim V_j = j\) and \(r < n = \dim V\). Then \(\mathfrak{g}\) acts on \(V/V_r\), and since the action of any \(x \in \mathfrak{g}\) on \(V\) is nilpotent, the same holds for the induced action on \(V/V_r\). Thus we apply \textbf{Proposition \ref{engelstheoremlemma}} again to get a common eigenvector \(w\) with eigenvalue zero, let \(w'\) be a lift of \(w\) to \(V\), then defining \(V_{r+1} = V_r \oplus \operatorname{span}\set{w'}\) we are done by induction.
            \end{proof}
        \end{theorem}

        \begin{theorem}[Engel]\label{engel}
            If every element of \(\mathfrak{g}\) is ad-nilpotent, then \(\mathfrak{g}\) is nilpotent.

            \begin{proof}
                By \textbf{Theorem \ref{engel2}} there is a basis for \(\mathfrak{g}\) such that the matrix of every element of \(\operatorname{ad}(\mathfrak{g})\) is strictly upper triangular. Since strictly upper triangular matrices in \(\mathfrak{gl}(\mathfrak{g})\) are a nilpotent lie algebra so is their subalgebra \(\operatorname{ad}(\mathfrak{g})\). This proves that \(\operatorname{ad}(\mathfrak{g}) \cong \mathfrak{g}/\mathfrak{z}(\mathfrak{g})\) is nilpotent. To see this implies \(\mathfrak{g}\) is nilpotent note that \(\mathfrak{g}/\mathfrak{z}(\mathfrak{g})\) being nilpotent implies that for some \(n > 0\) we have the central series term \(\mathfrak{g}^{n} \subset \mathfrak{z}(\mathfrak{g})\), then since \([\mathfrak{g},\mathfrak{z}(\mathfrak{g})] = 0\) we find that the central series term \(\mathfrak{g}^{n+1} = 0\) so that \(\mathfrak{g}\) is nilpotent.
            \end{proof}
        \end{theorem}
\end{document}